% Cote 119: offprints, reports, courses, CV and letters (typed and printed part): archive cote 119, site page 293 (PDF page 294).
% typescript. Draft compared: none.
% AI-transcribed by Claude Opus 5.5 (claude-opus-5-5, adaptive thinking, effort high), Message Batches API, batch r119_full (msgbatch_01DwZNm5Y7YsKBWyAN4sFs8w), 2026-11-13; not yet checked by eye.
\scanpage{119/293}
\mspage{- 6 -}

3.4. Le groupe $F_4$

Ici encore, on dispose d'une interprétation concrète de la cohomologie :
les éléments de $H^1(k, F_4)$ correspondent aux classes d'\uline{algèbres de Jordan}
\uline{simples exceptionnelles} de dimension 27 sur k.Malheureusement, on
est loin de savoir classer de telles algèbres, malgré les nombreux
résultats déjà obtenus par Albert, Jacobson, Tits, Springer, McCrimmon,
Racine, Petersson... Ces résultats suggèrent que les éléments de
$H^1(k, F_4)$ pourraient être caractérisés par deux types d'invariants :
(3.4.1 - ``\uline{invariant} mod 2'') La classe de la forme bilinéaire `` trace \hand{''}
associée à l'algèbre de Jordan, cette classe étant elle-même déterminée
par le couple d'une 3-\uline{forme de Pfister} et d'une 5-\uline{forme de Pfister} divisible
par la première. Du point de vue cohomologique, cela signifierait un
élément décomposable $x_3 \hand{$\in$} H^3(k, Z/2Z)$ (obtenu par (3.2.1) grâce à la
représentation irréductible \hand{$\rho$} de dimension 2\unsure{6} de $F_4$), et un élément
$x_5$ de $H^5(k, Z/2Z)$ de la forme $x_5 = x_3yz$ avec $y,z \hand{$\in$} H^1(k, Z/2Z)$.
(3.4.2 - `` \uline{invariant} mod 3 '') Un élément de $H^3(k, Z/3Z)$, dont je n'ai
qu'une définition conjecturale, basée sur la `` première construction
de Tits '' (on suppose ici que la caractéristique de k est \hand{$\neq$} 3).

Pour le moment, le seul cas qui puisse être traité complètement est
celui des algèbres de Jordan dites `` réduites '' (celles où l'invariant
mod 3 est 0): on sait, d'après un théorème de Springer, que l'invariant
mod 2 (i.e. la forme trace) détermine alors l'algèbre de Jordan à isomorphisme
près.

3.5. Le groupe $E_8$

Lorsque k est un corps de nombres, la structure de $H^1(k, E_8)$ vient
d'être déterminée par Chernousov et Premet : le principe de Hasse est
valable, ce qui entraîne par exemple que le nombre d'éléments de
$H^1(k, E_8)$ est $3^r$ , où r est le nombre de places réelles de k.La démons-%
tration de ce résultat a fait l'objet d'une série d'exposés dans le séminaire
commun avec la chaire de Théorie des Groupes.

Lorsque k est un corps quelconque (ou même, par exemple, le corps
Q(T)), on sait fort peu de choses sur $H^1(k, E_8)$. Les résultats généraux de
Grothendieck (\uline{sém}.\uline{Chevalley}, 1958) et de Bruhat-Tits (\uline{J}.\uline{Fac}.\uline{Sci}.\uline{Tokyo}
34, 1987) suggèrent qu'un élément\hand{|}de cet ensemble pourrait avoir comme
invariants des classe\hand{s} de cohomologie (de dimension \hand{$\geqslant$} 3) mod 2, mod 3
et mod 5 (car 2,3,5 sont les \uline{nombres premiers de torsion} de $E_8$ ,
cf. A.Borel, \uline{Oe}.II, p.776). J'ignore comment ces invariants pourraient être

\mspage{\hand{293}}
